\documentclass[12pt,a4paper]{article}
\usepackage{fontspec}\usepackage{polyglossia}\setdefaultlanguage{ukrainian}
\setmainfont{NewComputerModern10}\setsansfont{Noto Sans}\setmonofont{Noto Sans Mono}[Scale=0.82]
\usepackage{amsmath,amssymb,booktabs,caption,enumitem,float,geometry,graphicx,hyperref,microtype,titlesec,xcolor}
\geometry{left=2cm,right=2cm,top=2cm,bottom=2cm}\setlength{\parindent}{1.25cm}\linespread{1.08}
\titleformat{\section}{\normalfont\fontsize{14}{16}\selectfont\bfseries}{\thesection}{0.6em}{}\titlespacing*{\section}{0pt}{1.4em}{.65em}
\captionsetup{font=small,labelfont=bf,labelsep=endash}\usepackage[newfloat]{minted}\SetupFloatingEnvironment{listing}{name=Лістинг}\setminted{fontsize=\footnotesize,linenos,breaklines,autogobble,frame=single,framesep=2mm,bgcolor=gray!5,tabsize=4}
\hypersetup{unicode,colorlinks=true,linkcolor=black,pdftitle={ІО-41, Давидчук А. М., Лабораторна робота №4},pdfauthor={Давидчук Артем}}
\newcommand{\resultDir}{../results}\newcommand{\studentGroup}{ІО-41}\newcommand{\studentName}{Давидчук А. М.}\newcommand{\reportYear}{2026}
\begin{document}\pagenumbering{gobble}\begin{titlepage}\begin{center}\includegraphics[width=.7\textwidth]{../../../../../../templates/letterhead.png}\end{center}\begin{center}\large Національний технічний університет України\\ «Київський політехнічний інститут імені Ігоря Сікорського»\\[1em]Факультет інформатики та обчислювальної техніки\\Кафедра обчислювальної техніки\end{center}\vfill\begin{center}\textbf{\LARGE Методи та технології штучного інтелекту}\\[2em]\textbf{\Large Лабораторна робота №4}\\[.7em]«Моделювання функції двох змінних на основі нейронних мереж»\end{center}\vfill\begin{flushright}\begin{minipage}{.48\textwidth}Виконав: студент 3 курсу ФІОТ,\\групи \studentGroup\\\textit{\studentName}\\Перевірив: \textit{Польшакова О. М.}\end{minipage}\end{flushright}\vfill\begin{center}Київ --- \reportYear\end{center}\end{titlepage}\pagenumbering{arabic}

\section{Тема та мета роботи}
\noindent\textbf{Тема:} моделювання функції двох змінних нейронною мережею.\\\noindent\textbf{Мета:} дослідити вплив архітектури та кількості нейронів на якість моделювання функції з двома входами й одним виходом.

\section{Варіант і постановка задачі}
Для варіанта 6 задано
\begin{equation}y=\sin(|x|)+\cos(|x|),\qquad z=x\sin(y).\end{equation}
У цій роботі перша формула задає залежність допоміжної змінної $y$ від $x$, а друга визначає вихідну величину $z$. Для навчання мережі використовуються пари числових входів $(x,y)$ і відповідний еталонний вихід $z$. Потрібно навчити та порівняти шість мереж: feed-forward з одним шаром на 10 і 20 нейронів; cascade-forward з одним шаром на 20 нейронів та двома шарами по 10; Elman з одним шаром на 15 нейронів та трьома шарами по 5 нейронів.

Таке порівняння дає змогу перевірити, чи покращується наближення при збільшенні кількості нейронів, при додаванні прихованого шару та при використанні контекстних зв'язків Elman.

\section{Теоретичні відомості}
Штучна нейронна мережа складається з нейронів, ваг, зміщень та активаційних функцій. Для нейрона обчислюється $s=\sum_iw_ix_i+b$, після чого застосовується функція активації. Ваги визначають внесок кожного входу, а зміщення дозволяє змістити поріг спрацьовування нейрона.

У прихованих шарах використано гіперболічний тангенс:
\begin{equation}\tanh(s)=\frac{e^s-e^{-s}}{e^s+e^{-s}},\end{equation}
а на виході --- лінійну функцію, оскільки задача є регресійною і вихід $z$ не є класом. Навчання виконується методом зворотного поширення помилки з функцією втрат MSE:
\begin{equation}MSE=\frac{1}{n}\sum_{i=1}^{n}(z_i-\hat z_i)^2.\end{equation}

Feed-forward мережа передає сигнал лише вперед, тому кожен шар отримує результат попереднього шару. Cascade-forward додатково з'єднує початкові входи з наступними шарами. Це спрощує проходження інформації від входу та може допомогти для функцій із прямою залежністю від входів.

Elman-мережа має контекстний зв'язок: прихований стан попереднього кроку подається назад як додаткова інформація. У цій роботі мережа проходить точки навчальної вибірки послідовно, а попередній прихований стан використовується як контекст. Для статичної функції це не обов'язково дає перевагу, але дозволяє дослідити принцип роботи рекурентної архітектури.

\section{Програмна реалізація}
Дані сформовано на рівномірній сітці значень $x\in[0;1]$ та $y\in[1;1{,}3818]$. Для кожної точки обчислено еталонне значення $z=x\sin(y)$. Усього отримано 625 пар входів і виходів. 80\% точок використано для навчання, 20\% --- для незалежного тестування. Розбиття виконується випадковою перестановкою з фіксованим seed, тому результат можна відтворити.

Під час навчання мережа виконує пряме поширення сигналу, знаходить різницю між прогнозом та еталоном, а потім коригує ваги градієнтним спуском. Навчання повторюється 350 епох. Після цього прогноз на тестових точках порівнюється з точним значенням функції. Приклад обчислення функції та конфігурацій:
\begin{listing}[H]\begin{minted}{python}
def function_z(x, y):
    return x * np.sin(y)

configurations = [
    ("feedforward_1x10", "feedforward", [10]),
    ("feedforward_1x20", "feedforward", [20]),
    ("cascade_1x20", "cascade", [20]),
    ("cascade_2x10", "cascade", [10, 10]),
    ("elman_1x15", "elman", [15]),
    ("elman_3x5", "elman", [5, 5, 5]),
]
\end{minted}
\caption{Функція та досліджувані конфігурації}
\end{listing}

Середня відносна похибка розрахована як
\begin{equation}\varepsilon=\frac{1}{n}\sum_{i=1}^{n}\frac{|z_i-\hat z_i|}{|z_i+1|}\cdot100\%.\end{equation}

\section{Результати дослідження}
\begin{table}[H]\centering\small\begin{tabular}{lrrr}\toprule Мережа & Шари & Нейрони & Похибка, \%\\\midrule Feed-forward 1x10&1&10&1,7920\\Feed-forward 1x20&1&20&2,8325\\Cascade-forward 1x20&1&20&1,4647\\Cascade-forward 2x10&2&20&2,1152\\Elman 1x15&1&15&1,6477\\Elman 3x5&3&15&3,1417\\\bottomrule\end{tabular}\caption{Порівняння нейронних мереж}\end{table}
\begin{figure}[H]\centering\includegraphics[width=.9\textwidth]{\resultDir/training_loss.png}\caption{Залежність MSE від епохи навчання}\end{figure}
Найменшу тестову похибку отримано для cascade-forward з одним внутрішнім шаром і 20 нейронами: $1{,}4647\%$. Серед feed-forward кращою виявилась мережа з 10 нейронами. Для Elman-конфігурації збільшення кількості шарів у цьому експерименті погіршило результат.
\begin{figure}[H]\centering\begin{minipage}{.48\textwidth}\includegraphics[width=\textwidth]{\resultDir/surface_cascade_1x20.png}\end{minipage}\hfill\begin{minipage}{.48\textwidth}\includegraphics[width=\textwidth]{\resultDir/surface_feedforward_1x10.png}\end{minipage}\caption{Приклади поверхонь нейромережевих моделей}\end{figure}

На графіку навчання всі криві загалом зменшуються, тобто мережі вчаться відтворювати задану залежність. Найшвидше зменшення помилки не завжди означає найкращий тестовий результат, оскільки важлива здатність узагальнювати дані, яких мережа безпосередньо не бачила під час навчання.

Feed-forward 1x10 показала похибку $1{,}7920\%$. Збільшення числа нейронів до 20 у цій реалізації дало $2{,}8325\%$, тому саме збільшення ширини мережі не гарантувало покращення. Cascade-forward 1x20 дала найменшу похибку $1{,}4647\%$: прямі зв'язки входів із вихідною частиною виявилися корисними для цієї гладкої функції. Друга cascade-forward конфігурація з двома шарами мала похибку $2{,}1152\%$.

Elman 1x15 отримала $1{,}6477\%$, що близько до найкращого результату. Elman 3x5 дала $3{,}1417\%$. Це показує, що додавання контекстних шарів і збільшення глибини потрібно оцінювати експериментально, оскільки складніша мережа може навчатися важче або гірше узагальнювати.

\section{Висновки}
Побудовано шість нейронних мереж для наближення функції двох змінних. Усі моделі навчилися відтворювати загальну форму функції, але якість залежала від структури мережі. Найменшу середню відносну похибку показала cascade-forward мережа з одним прихованим шаром та 20 нейронами. Отже, збільшення кількості нейронів або шарів не гарантує покращення: результат залежить від архітектури та параметрів навчання.

\section{Контрольні питання}
\begin{enumerate}[leftmargin=*]
\item \textit{Що таке штучна нейронна мережа і за аналогією з чим вона створена?}\\
Штучна нейронна мережа --- математична модель, яка складається зі з'єднаних обчислювальних елементів і натхненна будовою біологічної нервової системи.
\item \textit{З яких основних елементів складається штучний нейрон?}\\
Нейрон має входи, ваги, зміщення, суматор і функцію активації.
\item \textit{Що таке архітектура нейронної мережі?}\\
Архітектура описує кількість шарів, кількість нейронів у шарах та схему з'єднань між ними.
\item \textit{Які бувають типи нейронних мереж?}\\
Основні типи --- мережі прямого поширення, рекурентні, згорткові та мережі зі спеціальними пропусками зв'язків.
\item \textit{Що таке активаційна функція?}\\
Це нелінійне перетворення зваженої суми входів, яке дає змогу мережі моделювати складні залежності.
\item \textit{Які активаційні функції найчастіше використовують?}\\
Поширені функції: ReLU, sigmoid, гіперболічний тангенс і softmax.
\item \textit{Що таке навчання нейронної мережі?}\\
Це підбір ваг і зміщень за навчальними прикладами так, щоб зменшити помилку прогнозу.
\item \textit{Яка різниця між навчанням з учителем, без учителя та з підкріпленням?}\\
З учителем використовуються еталонні відповіді, без учителя мережа шукає структуру сама, а з підкріпленням навчається за винагородою за дії.
\item \textit{Що таке функція втрат?}\\
Функція втрат чисельно оцінює різницю між прогнозом мережі та еталонним значенням. У роботі використано MSE.
\item \textit{Як здійснюється корекція ваг?}\\
Ваги змінюються градієнтним спуском у напрямку зменшення функції втрат.
\item \textit{Що таке backpropagation?}\\
Це алгоритм обчислення градієнтів, який поширює помилку від вихідного шару до попередніх шарів.
\item \textit{Які властивості роблять нейронні мережі ефективними?}\\
Нелінійність, здатність узагальнювати, адаптивність і можливість автоматично вивчати залежності.
\item \textit{Де застосовують нейронні мережі?}\\
У розпізнаванні зображень і мовлення, прогнозуванні, керуванні, діагностиці та робототехніці.
\item \textit{Яка перевага нейромереж над класичними алгоритмами?}\\
Нейромережа може сама налаштувати параметри за даними та апроксимувати складну нелінійну функцію.
\item \textit{Які недоліки та обмеження мають нейронні мережі?}\\
Їм потрібні якісні дані, підбір гіперпараметрів і значні обчислення; результат часто складно інтерпретувати.
\end{enumerate}
\end{document}
